\section{Metodología}
    El análisis se realizó en un cuaderno de Google Colab cuyos segmentos de
    código se reproducen a lo largo de este informe en el mismo orden en que se
    ejecutan, puesto que cada paso depende del anterior, iniciando con la
    preparación del entorno, siguiendo con la descarga y verificación del
    conjunto, la depuración de duplicados, el tratamiento de los valores faltantes
    y la conversión de tipos, la caracterización de la variable objetivo y, por
    último, el examen de su relación con cada grupo de variables explicativas. El
    enfoque sigue la lógica del análisis exploratorio de datos de Tukey, mirar los
    datos antes de modelarlos para descubrir su estructura y sus anomalías
    \cite{tukey1977exploratory}, dentro del marco de fases de CRISP-DM
    \cite{chapman2000crisp}.

    \subsection{Entorno de trabajo y reproducibilidad}
        El cuaderno utiliza únicamente librerías preinstaladas en Colab,
        \textit{kagglehub} para descargar el conjunto de datos,
        \textit{pandas} para la carga y agregación de datos
        \cite{mckinney2010pandas}, \textit{matplotlib} y \textit{seaborn} para las
        figuras \cite{hunter2007matplotlib}, \cite{waskom2021seaborn},
        \textit{scipy} para la correlación de Spearman y \textit{scikit-learn} para
        el ajuste lineal exploratorio \cite{pedregosa2011scikit}. Las rutas se
        resuelven respecto de la raíz del repositorio, de modo que las figuras se
        guardan en \texttt{assets/images/data/}, que es de donde las toma este
        informe.

        \vspace{10pt}

        \lstinputlisting{\repo utils/codes/environment-setup/environment-setup.py}

        \vspace{10pt}

        Luego se fija la identidad gráfica, un azul para la variable objetivo y
        las series principales, un naranja para contrastes y un tono neutro para
        las descripciones del conjunto, y la semilla global \texttt{RANDOM\_STATE
        = 42}, que gobierna el único muestreo del cuaderno.

        \vspace{10pt}

        \lstinputlisting{\repo utils/codes/graphic-identity/graphic-identity.py}

    \subsection{Descarga y carga del conjunto}
        El periodo 2022-4 se extrajo del Portal Nacional de Datos Abiertos
        mediante su API, con un filtro por periodo y un orden fijo por
        identificador, y se publicó en Kaggle como el dataset \textit{Saber 11
        Colombia - Resultados 2022-4} \cite{rubiano2026kaggle}, con la misma
        licencia y atribución al ICFES. La copia conserva todas las filas y
        columnas tal como las entrega el portal; el único cambio es el formato,
        de CSV (560 MB) a Parquet (31,5 MB), que además guarda el tipo de cada
        columna: las categóricas como \texttt{category}, los puntajes como
        \texttt{float32} y los identificadores como cadena. Publicar una copia
        congelada hace el análisis reproducible aunque el portal actualice sus
        datos, y reduce la carga del cuaderno a unos segundos. La celda descarga
        únicamente el archivo Parquet con \textit{kagglehub}, que no requiere
        cuenta de Kaggle para un dataset público, y lo lee con \textit{pandas}.

        \vspace{10pt}

        \lstinputlisting{\repo utils/codes/dataset-download/dataset-download.py}

        \vspace{10pt}

        La estructura del conjunto se verifica columna por columna, tipo,
        cantidad de valores no nulos, porcentaje de faltantes, cardinalidad y un
        valor de ejemplo, que es la base del diccionario de variables de la
        Sección~III.

        \vspace{10pt}

        \lstinputlisting{\repo utils/codes/dataset-structure/dataset-structure.py}

    \subsection{Calidad de los datos}
        Se cuentan los registros duplicados y se comprueba si el identificador
        \texttt{estu\_consecutivo} es único; los duplicados exactos se eliminan y
        el resto del análisis se realiza sobre el conjunto depurado. Se
        identifican además las columnas constantes y las de muy poca variación,
        que no aportan información a un modelo, y se tabula la cardinalidad de
        cada variable categórica, que determina su codificación en la fase 2.

        \vspace{10pt}

        \lstinputlisting{\repo utils/codes/data-quality/data-quality.py}

        \vspace{10pt}

        Los valores faltantes se cuantifican por columna y se examina si los del
        bloque \texttt{fami\_*} ocurren juntos, lo que indicaría estudiantes que
        no diligenciaron el cuestionario, o de forma independiente; también se
        compara el puntaje de quienes no tienen cuestionario con el del resto,
        para saber si la ausencia es informativa.

        \vspace{10pt}

        \lstinputlisting{\repo utils/codes/missing-values/missing-values.py}

        \vspace{10pt}

        Por último se convierte la fecha de nacimiento en la edad al momento de
        la prueba, tomando como referencia la fecha de aplicación del calendario A
        de 2022, y se fija el orden natural de las variables ordinales (estrato,
        educación de los padres, personas y cuartos del hogar, nivel de inglés),
        que \texttt{category} no conoce por sí misma.

        \vspace{10pt}

        \lstinputlisting{\repo utils/codes/type-conversion/type-conversion.py}

    \subsection{Variable objetivo}
        La distribución de \texttt{punt\_global} se caracteriza con sus
        estadísticas descriptivas, su asimetría y curtosis, un histograma y un
        diagrama de caja. Luego se verifica la fórmula con la que el ICFES
        construye el puntaje global a partir de las cinco pruebas
        \cite{icfes2022guia},
        \begin{equation}
            \label{global_formula}
            \begin{aligned}
                G = \operatorname{round}\Bigl(\tfrac{5}{13}\bigl(
                    & 3\,LC + 3\,MAT \\
                    & + 3\,SOC + 3\,CN + 1\,ING \bigr)\Bigr),
            \end{aligned}
        \end{equation}
        donde $LC$, $MAT$, $SOC$, $CN$ e $ING$ son los puntajes de lectura
        crítica, matemáticas, sociales y ciudadanas, ciencias naturales e inglés.
        Si la fórmula se cumple sobre los datos, los puntajes por área son
        componentes de la respuesta y no variables explicativas, y su inclusión en
        un modelo constituiría una fuga de información \cite{kaufman2012leakage}.

        \vspace{10pt}

        \lstinputlisting{\repo utils/codes/target-distribution/target-distribution.py}

        \vspace{10pt}

        \lstinputlisting{\repo utils/codes/target-formula/target-formula.py}

    \subsection{Análisis bivariado por grupo de variables}
        Para cada variable categórica se calcula el tamaño, la media, la mediana y
        la desviación del puntaje global por categoría, y se grafica con diagramas
        de caja sin valores atípicos, de modo que la comparación entre grupos no
        quede dominada por los extremos. Dos funciones auxiliares,
        \texttt{resumen\_por} y \texttt{caja\_por}, encapsulan ese patrón y se
        reutilizan en las secciones de estudiante, hogar y colegio. La edad se
        agrupa en tramos, porque su relación con el puntaje no tiene por qué ser
        lineal.

        \vspace{10pt}

        \lstinputlisting{\repo utils/codes/student-variables/student-variables.py}

        \vspace{10pt}

        \lstinputlisting{\repo utils/codes/family-variables/family-variables.py}

        \vspace{10pt}

        Para las variables binarias del hogar se compara la media del puntaje
        entre quienes tienen y no tienen cada bien, y se reporta la diferencia en
        puntos.

        \vspace{10pt}

        \lstinputlisting{\repo utils/codes/family-binary/family-binary.py}

        \vspace{10pt}

        \lstinputlisting{\repo utils/codes/school-variables/school-variables.py}

        \vspace{10pt}

        Como naturaleza y área se cruzan (la mayoría de colegios rurales son
        oficiales), se tabula el puntaje medio en las cuatro combinaciones.

        \vspace{10pt}

        \lstinputlisting{\repo utils/codes/school-crosstab/school-crosstab.py}

        \vspace{10pt}

        Para la geografía se calcula el puntaje medio por departamento del
        colegio, se ordena y se contrasta con la media nacional.

        \vspace{10pt}

        \lstinputlisting{\repo utils/codes/geography/geography.py}

    \subsection{Correlaciones y ajuste lineal exploratorio}
        Las variables ordinales se codifican con enteros que respetan su orden,
        estrato de 0 a 6 y nivel educativo de 0 a 9 (las categorías \textit{No
        sabe} y \textit{No Aplica} se tratan como faltantes), las binarias con 0 y
        1, y se calcula la correlación de Spearman \cite{spearman1904proof},
        \begin{equation}
            \label{spearman}
            \rho = 1 - \frac{6 \sum_{i=1}^{n} d_i^2}{n\,(n^2 - 1)},
        \end{equation}
        donde $d_i$ es la diferencia entre los rangos de las dos variables para la
        observación $i$. Se eligió Spearman en lugar de Pearson porque mide
        asociaciones monótonas, no solo lineales, y es apropiada para variables
        ordinales. La matriz incluye los puntajes por área para documentar el
        contraste entre una variable explicativa legítima y una que forma parte de
        la respuesta.

        \vspace{10pt}

        \lstinputlisting{\repo utils/codes/correlations/correlations.py}

        \vspace{10pt}

        Finalmente se ajusta una regresión lineal exploratoria sobre las
        variables codificadas, excluidos los puntajes por área, con una muestra
        aleatoria de 200.000 filas completas y una partición 80/20. No es el
        modelo de referencia de la fase 2, que se construirá con la limpieza
        completa y la partición definitiva, sino una cota rápida de cuánta
        varianza del puntaje explican las variables candidatas en conjunto,
        medida con el coeficiente de determinación $R^2$ y el error absoluto medio
        (MAE).

        \vspace{10pt}

        \lstinputlisting{\repo utils/codes/exploratory-fit/exploratory-fit.py}
